 \section*{Chapitre \ref{ch:s1:intraday-patterns} --- Régularités intrajournalières}

\begin{solution}{exo:s1:intraday-patterns:1}
$16~\%/\sqrt{20} = 3.6~\%$ par an pour le rendement total~; la part overnight a une volatilité de $16~\% \times \sqrt{0.25} = 8~\%$, donc $8~\%/\sqrt{20} = 1.8~\%$ par an.
\end{solution}

\begin{solution}{exo:s1:intraday-patterns:2}
Un quart~: $0.25 \times 7~\% = 1.75~\%$ par an overnight, 5.25~\% en journée.
\end{solution}

\begin{solution}{exo:s1:intraday-patterns:3}
$0.6 \times 10 = 6$ points de base par écart type de déséquilibre, avant coûts~; six points de base chaque jour donneraient $6 \times 252/100 = 15.1~\%$ par an. Le
La stratégie du modèle rapporte 8.8~\%, car les déséquilibres varient en taille, elle prend une position fixe quel que soit le déséquilibre, et elle supporte un point de base par jour.
\end{solution}

\begin{solution}{exo:s1:intraday-patterns:4}
$0.05 \times 2~\% = 0.10~\%$, soit dix points de base de pression acheteuse dans la dernière demi-heure.
\end{solution}

\begin{solution}{exo:s1:intraday-patterns:5}
Le delta d'un book à gamma court baisse quand le cours monte (il devient plus court), si bien que son hedger achète après les hausses et vend après les baisses~: dans le sens du mouvement. Un book à gamma long
a un delta qui monte avec le cours, et son hedger vend dans les hausses et achète dans les baisses~: à contre-sens du mouvement.
\end{solution}

\begin{solution}{exo:s1:intraday-patterns:6}
Une moyenne de rendement sur vingt ans a une erreur type de 3.6 points de pourcentage par an, et sa part overnight 1.8 points de pourcentage~: le quart injecté dans le modèle a été mesuré à 54~\%.
Le résultat de Bondarenko et Muravyev concentre le rendement sur quatre heures, avec un ratio de Sharpe de 1.6, ce qui le rend détectable, contrairement à une découpe
d'une moyenne annuelle bruitée, qui ne l'est pas.
\end{solution}

\begin{solution}{exo:s1:intraday-patterns:7}
La contre-tendance ne dépend que du déséquilibre de clôture et de son \omterm{def:s1:short-term-reversal:overnight}{retournement overnight}, qui ne dépendent pas de la couverture des dealers~: 1.09 et 1.08. La
transaction n'existe que si la couverture fait suivre la journée à la dernière demi-heure~: $-0.77$ sans elle, 1.05 avec elle.
\end{solution}

\begin{solution}{exo:s1:intraday-patterns:8}
Une moyenne d'échantillon des \omterm{def:s1:intraday-patterns:returns}{rendements overnight} moins les rendements de jour est bruitée (le modèle mesure 54~\% en overnight quand un quart est injecté), et ne détenir que l'overnight
divise par deux le temps investi et double le nombre de transactions. Tester la découpe sur un long échantillon continu sur vingt-quatre heures, avec coûts, et se demander pourquoi la prime serait gagnée
la nuit.
\end{solution}

\begin{solution}{pb:s1:intraday-patterns:1}
\begin{enumerate}
\item De la clôture à l'ouverture suivante~; de l'ouverture à la clôture.
\item Le rendement moyen total de l'E-mini sur quatre heures autour de l'ouverture européenne~; les profits de retournement et de momentum réalisés overnight, avec un retournement entre périodes.
\item Un quart injecté (1.75~\% par an)~; 54~\% mesurés (4.4~\% en overnight, 3.7~\% en journée).
\item Erreur d'échantillonnage~: une moyenne sur vingt ans est trop bruitée pour être découpée.
\item Le \omterm{def:s1:intraday-patterns:smile}{smile du volume}, en forme de U au cours de la journée~; 12.1~\% dans chacune des première et dernière demi-heures, 4.9~\% à midi.
\item Avant 15 h 50, ordres à la clôture~; à partir de 15 h 50, données de déséquilibre~; heure limite MOC 15 h 55, LOC 15 h 58~; croisement à 16 h.
\item Les \omterm{def:s1:index-rebalancing:fund}{fonds indiciels}, les fonds valorisés à la clôture, les dealers qui valorisent leurs books~: c'est là que se fixe le prix de l'indice de référence.
\item 8.8~\% par an, avec un ratio de Sharpe de 1.09.
\item Négocier près de la clôture pour rééquilibrer l'exposition après le mouvement de la journée~; dans le sens du mouvement lorsque le gamma est court.
\item Les 30 dernières minutes suivent le reste de la journée sur plus de soixante contrats à terme, 1974--2020, avec un retournement dans les jours suivants, lié à la couverture gamma.
\item Pente $-0.006$ et transaction perdante ($-0.77$) sans couverture~; 0.044 ($t = 9.2$) et ratio de Sharpe de 1.05 avec couverture.
\item La couverture opère à contre-sens du mouvement et la régularité s'inverse.
\item \textbf{Résultat nommé.} 54~\% du rendement cumulé du modèle est mesuré en overnight là où un quart est injecté, une erreur d'échantillonnage d'une moyenne sur vingt ans~;
la pente de la dernière demi-heure sur le reste de la journée vaut $-0.006$ ($t = -1.2$) sans flux de couverture et 0.044 ($t = 9.2$) avec lui.
\item La fenêtre overnight est une prime de risque~; les transactions sur la dernière demi-heure et sur déséquilibre sont des transactions de liquidité.
\item À partir des positions ouvertes sur options et du côté probable des dealers, agrégés en gamma net par strike.
\item Fourchettes et frais sur les allers-retours quotidiens~; frais d'enchère.
\item L'ensemble des données de contrats à terme sur vingt-quatre heures, leurs coûts et la stabilité du résultat d'une période à l'autre.
\item La fenêtre overnight sur les contrats à terme sur indices.
\item Les jours de reconstitution concentrent le flux indiciel dans l'enchère de clôture, là où la transaction sur déséquilibre et la fourniture de liquidité du chapitre~10 se rencontrent.
\item C'est là que se fixent les prix des indices de référence, si bien que ceux qui doivent négocier à ces prix paient ceux qui peuvent attendre.
\end{enumerate}
\end{solution}

\begin{solution}{iq:s1:intraday-patterns:1}
À l'ouverture, les nouvelles overnight et les ordres sont négociés~; à la clôture, les prix des indices de référence se fixent, et les fonds, les suiveurs d'indices et les hedgers y négocient~; à midi, on ne trouve
ni l'un ni l'autre.
\end{solution}

\begin{solution}{iq:s1:intraday-patterns:2}
Le rendement de la dernière demi-heure tend à suivre le rendement de la journée jusqu'à ce moment. La couverture des dealers à gamma court et des ETF à effet de levier, qui doivent négocier dans le
sens du mouvement près de la clôture, et la négociation informée tardive.
\end{solution}

\begin{solution}{iq:s1:intraday-patterns:3}
Prendre la contrepartie d'un déséquilibre important annoncé dans l'enchère, dimensionné à la profondeur de l'enchère, et sortir à l'ouverture suivante ou pendant la journée suivante, en évitant
les actions où le déséquilibre peut traduire une nouvelle information.
\end{solution}

\begin{solution}{iq:s1:intraday-patterns:4}
Utiliser les contrats à terme sur vingt-quatre heures, et non les cours d'ouverture et de clôture~; calculer les rendements par heure~; tenir compte des coûts~; tester la stabilité selon les périodes et les marchés~; et comparer
la taille de l'effet à l'erreur d'échantillonnage de la découpe.
\end{solution}

\begin{solution}{iq:s1:intraday-patterns:5}
Risque de gap à l'ouverture, sans possibilité de sortir~; nouvelles concentrées overnight~; le coût de deux transactions par jour~; et, si la prime rémunère
le risque overnight, ses pertes surviennent les mauvaises nuits.
\end{solution}

\begin{solution}{iq:s1:intraday-patterns:6}
Le delta d'un book de gamma $\Gamma$ varie de $\Gamma\,\Delta S$ après un mouvement $\Delta S$~; un dealer à gamma court ($\Gamma < 0$) doit acheter $|\Gamma|\Delta S$ pour rester
couvert après une hausse. Sous un impact linéaire, le cours se déplace de $\lambda |\Gamma| \Delta S$, dans le sens du mouvement initial~: une rétroaction proportionnelle au
gamma des dealers.
\end{solution}
